\section{不足与展望}

当前 DeepRTE 的实现是针对给定几何形状学习解算子。与大多数主流算子学习框架类似，当几何形状发生改变时，原则上需要重新训练模型。不过，可以采用若干策略缓解这一限制：

一是可在一个包含所有目标几何的规则区域（例如本文实验中的正方形区域）上训练模型，并将实际几何的边界条件适当延拓到该区域。当在该包含域内部应用到新的几何形状时，可以直接使用已有模型而无需完全重新训练。

二是引入以神经网络为基础的几何编码器与解码器，将不同几何映射到统一的潜在域，再从潜在域映回物理域。DeepRTE 模型在该潜在域上进行训练；遇到新几何时，只需利用少量数据，通过迁移学习微调编码器、解码器及相关模块即可。

此外，当前DeepRTE主要针对于无源项的辐射传输方程，对于源项我们可以采取叠加原理进行处理，然而由于线性性的缺失，加入源项的框架可能无法实现零样本泛化能力。未来可以尝试将源项作为额外的输入参数引入模型，或者设计专门处理源项的网络模块，以提升模型在有源项情况下的表现。

同时DeepRTE目前只处理了一个能级的辐射传输问题，未来可以考虑将多能级耦合的辐射传输方程纳入模型设计中，以应对更复杂的物理场景。这个处理应当是相对直接的，因为多能级耦合的RTE仍然保持线性结构，且各能级之间的耦合关系可以通过适当的网络设计进行建模。

\section{本章总结}

本章系统构建了辐射输运方程（RTE）的物理信息驱动神经算子求解框架DeepRTE。从RTE解的内在数学结构出发，我们对输运过程进行了可解释的物理分解，将原问题改写为等价的不动点迭代形式，为后续结合神经网络奠定了严谨的理论基础。

针对分解后关键算子难以获得解析表达的挑战，我们提出采用神经网络对衰减/传播映射、散射映射等核心组件进行参数化近似，并将其无缝嵌入到不动点迭代流程中，形成了端到端的残差神经算子架构以求解RTE。该框架的核心创新在于同时利用了物理先验（来自结构分解与格林函数理论）与数据驱动学习的强大表达能力，因而能够更有效地刻画辐射传输解中的结构。

作为一个端到端的神经算子学习框架，DeepRTE展现出显著的灵活性、高效性与泛化性，其核心优势可归纳为以下几个方面：

\paragraph{端到端的神经算子学习范式}
DeepRTE本质上是一个学习无限维函数空间之间映射的神经算子，而非传统的固定网格上的离散值拟合。它直接以介质参数（$\mu_t, \mu_s$）、边界条件（$I_-$）和散射核（$p$）等函数作为输入，直接输出全相空间的辐射强度场$I(\boldsymbol{r},\boldsymbol{\Omega})$，实现了从“物理输入”到“物理输出”的端到端学习，避免了传统方法中复杂的人工设计求解流程。

\paragraph{物理与数据融合的双驱动建模}
DeepRTE并非纯黑盒的数据驱动模型，我们在本章设计中深度融合了物理先验知识。网络架构严格遵循RTE的解结构，衰减模块与散射模块的设计分别对应于输运过程中的不同物理机制，格林函数积分的形式内嵌于网络的前向传播中，使得网络在小样本下即可收敛，且学习到的解具有物理一致性。

\paragraph{网格解耦与几何灵活性}
这是DeepRTE最核心的优势之一。输入网格点$\{\boldsymbol{r}^{\text{mesh}}_i\}$仅用于刻画介质参数$\mu_t$和$\mu_s$的空间分布，而与辐射场$I$的求解位置完全解耦。这带来了三重便利：
\begin{enumerate}
    \item \textbf{求解位置自由}：网格点不必与求解点重合，可在任意感兴趣的位置（如非结构化点云、局部加密区域）查询辐射强度，不受求解区域几何形状的限制；
    \item \textbf{介质表示高效}：对于$\mu_t$变化平缓或分段常数的区域，可以使用非常稀疏的网格来精确表示介质属性，从而大幅降低计算量；
    \item \textbf{非结构化网格支持}：天然适配非结构化网格与复杂几何边界，能够灵活适应实际应用中的复杂介质分布。
\end{enumerate}

\paragraph{高效的推理与计算可扩展性}
在推理阶段，DeepRTE无需像传统扫掠法那样进行多轮迭代，仅需一次前向传播即可得到全相空间的解。同时，架构天然支持局部化策略：在大规模三维问题中，可以仅选取特征线邻域内的网格点子集参与注意力计算（如带状区域、锥体区域等），从而在保持精度的同时显著提升推理速度。网络的模块化设计也使其易于通过增加网络深度或宽度来提升表达能力，具有良好的可扩展性。

